\documentclass[11pt]{article}
\usepackage[a4paper,margin=1.1in]{geometry}
\usepackage[T1]{fontenc}
\usepackage{lmodern}
\usepackage{amsmath,amssymb,amsthm}
\usepackage{booktabs}
\usepackage{array}
\usepackage{microtype}
\usepackage[hidelinks]{hyperref}
\usepackage{natbib}
\usepackage{setspace}
\usepackage{enumitem}

\newtheorem{theorem}{Theorem}[section]
\newtheorem{proposition}[theorem]{Proposition}
\newtheorem{lemma}[theorem]{Lemma}
\newtheorem{corollary}[theorem]{Corollary}
\theoremstyle{definition}
\newtheorem{definition}[theorem]{Definition}
\theoremstyle{remark}
\newtheorem{remark}[theorem]{Remark}

\newcommand{\leanunderscore}{\char`\_\hskip 0pt plus 0.4pt\relax}
\newcommand{\lean}[1]{\begingroup\ttfamily\small\let\_\leanunderscore #1\endgroup}
\newcommand{\formal}[1]{\par\noindent\begingroup\footnotesize\raggedright\emph{Formal statement:} \lean{#1}.\par\endgroup}
\newcommand{\Thorn}{\text{\th}}
\newcommand{\R}{\mathbb{R}}
\newcommand{\E}{\mathbb{E}}

\onehalfspacing
\tolerance=600 \emergencystretch=1.5em

\title{A Unique Stationary Equilibrium at Relative Risk Aversion Three:\\
How the Result Is Built}
\author{LeanEconomics Org [mostly Claude]\thanks{The Lean development described here is at
\url{https://github.com/LeanEconomics/LeanEconomics}. It was written mostly by Claude (Anthropic),
an AI coding model, working under the direction of Robert Kirkby (Victoria University of
Wellington). This note is a self-contained account of one result of the overlapping-generations
write-up in the same repository, \emph{Overlapping Generations in Lean: The Life-Cycle Household},
which it cites as the companion; the results it builds on are stated with their formal names and
not reproved. It is a project record, not a paper for submission.}}
\date{October 2026}

\begin{document}
\maketitle

\begin{abstract}
\noindent
In a life-cycle economy with uninsurable earnings risk, a borrowing constraint, sixty-period lives
and relative risk aversion three, at one rational calibration, there is exactly one stationary
equilibrium interest rate in $(-8\%,6\%]$, and it lies in $(4\%,6\%)$. The route is single crossing in
derivative form: excess
demand for capital is strictly increasing on the bracket because the firm's demand falls faster
than household supply can, and a certified lower bound on the slope of supply, built from rigorous
enclosures of the sixty consumption functions, sits at most a third of the way to the slope of
demand. The note states the results the construction rests on --- the stage properties of the
finite backward induction, the propensity floor, the cohort envelope, continuity of supply, the
interval Euler step --- and then explains how they are assembled: the derivative of consumption in
the rate from the Euler equality, the interval recursion that bounds it on grid intervals, the
frozen-transport argument that carries a bound from a rate to a cell of rates by continuation, and
the three hundred and seventy cells that cover $[2.3\%,6\%]$; below $2.3\%$ Lean theorems take
over. Every step but the arithmetic is a theorem; the
arithmetic is exact rational arithmetic rounded outward, evaluated by a program rather than
replayed in the kernel. The note closes with what that standing is and what would raise it.
\end{abstract}

\section{Introduction}

The companion write-up proves that the stationary equilibrium of a life-cycle Bewley economy is
unique at logarithmic utility, by Light's three theorems in the life-cycle setting: saving rises
with the interest rate at every state (Theorem 1), ordered policies give ordered capital supply
(Theorem 2), and a rising supply meets the firm's falling demand at most once (Theorem 3). Above
logarithmic utility Theorem 1 is false in fact --- at relative risk aversion five it fails pointwise
at the youngest ages in the top earnings states --- and no analytic route around it is known. What
this note records is the result that reaches $\gamma=3$ anyway, and the construction that reaches it.

The result is a statement about one economy at one calibration. Fix sixty-period lives, a
seven-state earnings chain, $\beta=24/25$, $\gamma=3$, and a Cobb--Douglas firm. Then on the
range $(-8\%,6\%]$ there is exactly one stationary equilibrium, and it lies in $(4\%,6\%)$. The
argument is single
crossing, but not in the form Theorem 1 supplies. Single crossing needs only that excess demand
$E(r)=S(r)-D(r)$, supply minus demand for capital, be strictly increasing; and since demand falls
steeply at the calibration, supply may fall too, as long as it falls more slowly. The construction
certifies a lower bound on the slope of supply, $S'(r)\ge-B_d$ with $B_d\le9.2$, against a demand
slope of magnitude at least $28.7$, everywhere on the bracket; and it certifies that supply is below
demand at $4\%$ and above it at $6\%$. Continuity of supply, which is a theorem, then gives an
equilibrium in the bracket, and strict monotonicity of $E$ gives at most one.

The word \emph{certifies} carries the standing of the result and the note is careful with it. Three
kinds of statement appear. Theorems checked by the Lean kernel, which are stated with their formal
names and not reproved: these are the companion's. Paper arguments, which are short and are given
in full here: the derivative of the Euler equality, the mean-value link from a derivative bound to a
secant bound, and a continuation argument. And arithmetic, which is exact rational arithmetic
rounded outward, carried out by a program; its inequalities are instances of the formalised
interval Euler step and of the recursions this note writes down, but the program that evaluates
them has not been replayed in the kernel. Section~\ref{sec:standing} tabulates which step is which.

The note is organised as a construction. Section~\ref{sec:economy} fixes the economy and states the
theorem. Section~\ref{sec:known} states the known results. Section~\ref{sec:slope} reduces
uniqueness to a bound on the slope of supply, Section~\ref{sec:deriv} derives the derivative of
consumption in the rate and the envelope that turns its failure into a slope bound,
Section~\ref{sec:encl} describes the enclosures of the consumption functions,
Section~\ref{sec:recursion} the interval recursion that bounds the derivative,
Section~\ref{sec:cell} the passage from a rate to a cell of rates, Section~\ref{sec:bracket} the
covering of the bracket, Section~\ref{sec:level} the endpoint certificate, and
Section~\ref{sec:proof} assembles the proof. Section~\ref{sec:standing} is about standing and what
remains.

\section{The economy and the statement}\label{sec:economy}

\subsection{Households}

A household lives $J=60$ periods, is born with no assets, and dies leaving nothing. Earnings
$y(z)$ follow a finite Markov chain with transition matrix $\pi$ on $n_z=7$ states; the gross
return on assets is $R=1+r$; cash on hand is $m=Ra+y(z)$; saving $a'$ is bounded below by a zero
borrowing limit and above by an artefactual cap $\bar a$; utility is CRRA, $u'(c)=c^{-\gamma}$, with
$\gamma=3$ and discount factor $\beta$. Periods are counted from the end of life. With $k$ periods
still to come after the present one, the continuation value is the $k$-th Bellman iterate from
zero,
\[
  V_k=T^k0,\qquad (Tv)(a,z)=\max_{a'\in[0,\bar a]}\Bigl\{u(m-a')+\beta\sum_{z'}\pi(z,z')\,v(a',z')\Bigr\},
\]
the stage-$k$ policy is $g_k=\operatorname{policyOf}V_k$ and the stage-$k$ consumption is
$c_k=m-g_k$ (\lean{stageValue}, \lean{stagePolicy}, \lean{stageConsumption}). Age $j$ of life is
stage $k=J-1-j$; in the last period the household consumes its cash on hand, $c_0=m$. The problem is
a finite backward induction, so every property of the stage policies is proved by induction on $k$
with an explicit base case, and nothing is a fixed point.

\subsection{Earnings}

The chain is the Rouwenhorst discretisation with persistence $\rho=9/10$ on seven states. Its
invariant law is binomial, $\nu=(1,6,15,20,15,6,1)/64$, and the earnings levels are the seven
values $(0.301,0.449,0.670,1,1.492,2.226,3.320)$ divided by their $\nu$-mean $1.1266$, so that mean
earnings are one:
\[
  y=(0.2672,\ 0.3986,\ 0.5947,\ 0.8876,\ 1.3244,\ 1.9759,\ 2.9470).
\]
All of these are rationals, and so is every quantity the certificate computes.

\subsection{The firm and the equilibrium}

Technology is Cobb--Douglas with capital share $\alpha=9/25$ and depreciation $\delta=2/25$. In
wage units --- capital per unit of the wage, the normalisation homotheticity allows --- the firm's
demand for capital at rate $r$ is
\[
  D(r)=\frac{\alpha}{(1-\alpha)(r+\delta)},\qquad
  D'(r)=-\frac{\alpha}{(1-\alpha)(r+\delta)^2},
\]
so $D(4\%)=4.6875$, $D(6\%)=4.0179$, and $|D'|$ falls from $39.06$ at $4\%$ to $28.70$ at $6\%$.

Every cohort is born with no assets and an earnings state drawn from $\nu$; cohorts have equal
weight. Write $(\mathcal S_kh)(a,z)=\sum_{z'}\pi(z,z')\,h(g_k(a,z),z')$ for the expectation of $h$
next period (\lean{stageStep}). The assets a household holds over the rest of its life, seen from its
current state, satisfy
\[
  A_0(a,z)=a,\qquad A_{k+1}(a,z)=a+(\mathcal S_{k+1}A_k)(a,z)
\]
(\lean{cohortAssets}), and capital supply per head is
\[
  S(r)=\frac1J\sum_z\nu(z)\,A_{J-1}(0,z)
\]
(\lean{olgCapital}): a finite sum of finite sums, with no stationary distribution anywhere. A
stationary equilibrium is a rate at which $S(r)=D(r)$.

\subsection{The statement}

\begin{theorem}[Uniqueness on the bracket at $\gamma=3$]\label{thm:main}
In the economy above, with $J=60$, $\beta=24/25$, $\gamma=3$, the seven-state chain, $\alpha=9/25$
and $\delta=2/25$, there is exactly one stationary equilibrium interest rate in $(-8\%,6\%]$, and
it lies in $(4\%,6\%)$. Specifically: $S$ is continuous, $S(4\%)<D(4\%)$ and $S(6\%)>D(6\%)$, so
an equilibrium exists in $(4\%,6\%)$; $E=S-D$ is strictly increasing on $[2.3\%,6\%]$, so there is
at most one there; and $S<D$ on $(-8\%,2.3\%]$ (Theorem~\ref{thm:outer}), so there is none there.
\end{theorem}

The theorem is silent above $6\%$. Theorem~\ref{thm:outer} excludes $[8.7\%,20\%]$ by proof term,
so $(6\%,8.7\%]$ is the only range below $20\%$ in which an equilibrium is not yet excluded;
closing it is machine time, not new argument, and Section~\ref{sec:standing} returns to it.

\section{The results the construction rests on}\label{sec:known}

Each result below is a theorem of the companion development, checked by the Lean kernel, and is
used here as a black box. The formal name follows each statement.

\paragraph{Stage properties.} At every stage $k$, every state and every asset level in the region:
consumption is positive (\lean{stageConsumption\_pos}); saving and consumption are both nondecreasing
in assets (\lean{stagePolicy\_mono}, \lean{stageConsumption\_mono}); consumption is concave in assets
(\lean{concaveOn\_stageConsumption}, Carroll--Kimball); and the two Euler inequalities hold, under the
cap,
\[
  \beta R\sum_{z'}\pi(z,z')\,u'\bigl(c_k(g_{k+1}(a,z),z')\bigr)\ \le\ u'\bigl(c_{k+1}(a,z)\bigr),
\]
at every state (\lean{stage\_euler\_le}), and the reverse inequality at every state where the
household saves (\lean{stage\_euler\_ge}). Where the household saves, the two together are the Euler
\emph{equality}.

\paragraph{The propensity floor.} Let $\Thorn=(\beta R)^{1/\gamma}$ and define the finite-horizon
minimal propensities to consume by $\kappa_0=1$, $\kappa_{k+1}=\kappa_kR/(\Thorn+\kappa_kR)$, so that
$1/\kappa_k=\sum_{i\le k}(\Thorn/R)^i$ (\lean{stageMPC}). Then $c_k(a,z)\ge\kappa_km$
(\lean{stageMPC\_mul\_le\_stageConsumption}), and, the sharper statement that the argument needs, the
propensity to consume out of assets is at least $\kappa_k$:
\begin{equation}\label{eq:mpcfloor}
  \kappa_kR\,(a_2-a_1)\ \le\ c_k(a_2,z)-c_k(a_1,z),\qquad a_1\le a_2,
\end{equation}
so saving is Lipschitz in assets with constant $(1-\kappa_k)R$
(\lean{stageMPC\_mul\_le\_stageConsumption\_sub}). A lower bound on $\kappa_k$ follows from any upper
bound $t\ge\Thorn$ as $1/\sum_{i\le k}(t/R)^i$ (\lean{stageMPC\_ge\_of\_patience\_le}).

\paragraph{The cohort slope.} With $G_0=1$ and $G_{k+1}=1+G_k(1-\kappa_{k+1})R$
(\lean{cohortSlope}), accumulated assets are Lipschitz in assets,
$A_k(a_2,z)-A_k(a_1,z)\le G_k(a_2-a_1)$ for $a_1\le a_2$ in the region
(\lean{cohortAssets\_sub\_le\_cohortSlope}).

\paragraph{The earnings marginal.} Since $\nu$ is invariant for $\pi$ and the cohort is born from
$\nu$, the $\nu$-weighted expectation of any function of the earnings state is unchanged by any
number of steps of the cohort recursion (\lean{sum\_mul\_incomeStep\_iterate}). An expectation of a
failure over the cohort is therefore bounded by $\sum_z\nu(z)\,\sup_a(\text{failure at }(a,z))$:
exact in earnings, worst-case only in assets.

\paragraph{Theorem 1 is an affine bound.} Since $a'=Ra+y(z)-c$ at either rate, for $r_1<r_2$,
\begin{equation}\label{eq:thm1}
  g_k(a,z;r_2)\ge g_k(a,z;r_1)\quad\Longleftrightarrow\quad
  c_k(a,z;r_2)-c_k(a,z;r_1)\le(r_2-r_1)\,a
\end{equation}
(\lean{stagePolicy\_le\_iff\_consumption\_le}). The failure of Theorem 1 at a state is the positive
part of the excess of the left side over the right.

\paragraph{The cohort envelope.} Couple the two economies on a common earnings path from the common
start $a_0=0$. If, at every age, the failure of Theorem 1 at the lower-rate state is at most $e_j$ in
$\nu$-expectation and the higher-rate saving is Lipschitz with constant $L_j$ in assets, then
\begin{equation}\label{eq:envelope}
  S(r_2)-S(r_1)\ \ge\ -\frac1J\sum_{j}f_j,\qquad f_0=0,\quad f_{j+1}=L_jf_j+e_j,
\end{equation}
which in cohort form, with the regions of the lower-rate economy and $G_k$ as the accumulated
Lipschitz constants, is \lean{olgCapital\_sub\_ge\_of\_failure}. The pointwise ordering of the
policies is not assumed; its failure is priced.

\paragraph{Continuity, existence, single crossing.} Capital supply is continuous in the rate
(\lean{continuousOn\_olgSupply}); if it lies below a continuous demand at the bottom of an interval
and above it at the top, the two meet (\lean{exists\_olgEquilibrium}); and a supply that never falls
meets a strictly falling demand at most once (\lean{olgEquilibriumRate\_unique}). The last is used
here in its elementary form --- a strictly increasing $E$ has at most one zero --- which needs no
theorem.

\paragraph{Reachable regions.} The artefactual cap must be slack where the household can be. A
family of asset sets, one per stage, each inside $[0,\bar a]$, each carried by the next stage's
policy into the previous set, and on each of which the policy stays strictly below the cap, is a
reachable family (\lean{StageRegions}); a cohort born with nothing stays in its family
(\lean{reachRegions}), and every stage property above holds on any reachable family. Below, the
family is the enclosure's own reach: the sets $[0,\rho_j]$ with $\rho_j$ the certified maximal
assets at age $j$, about $44$ at most, far below any cap the computation uses.

\paragraph{The interval Euler step.} Let $E_k(x)=\beta R\sum_{z'}\pi(z,z')\,x(z')^{-\gamma}$ and
$m=y_z+Ra$.
\begin{itemize}[leftmargin=1.5em,itemsep=0pt]
\item If $0<c_L\le m$ and $E_k(\ell)\le c_L^{-\gamma}$ for numbers $\ell(z')\le c_k(m-c_L,z')$, then
  $c_L\le c_{k+1}(a,z)$ (\lean{le\_stageConsumption\_of\_cert}).
\item If $0<c_H\le m$, $m-c_H$ lies in the region, and $c_H^{-\gamma}\le E_k(h)$ for numbers
  $h(z')\ge c_k(m-c_H,z')$, then $c_{k+1}(a,z)\le c_H$ (\lean{stageConsumption\_le\_of\_cert}).
\item If $E_k(\ell)\le m^{-\gamma}$ for $\ell(z')\le c_k(0,z')$, then $c_{k+1}(a,z)=m$: the
  household is at the corner (\lean{stageConsumption\_eq\_resources\_of\_cert}).
\item Between grid points a concave nondecreasing function lies above the chord of lower bounds at
  two points and below the left secant extended forward and the right secant extended backward
  (\lean{le\_of\_chord}, \lean{le\_of\_left\_secant}, \lean{le\_of\_right\_secant}).
\end{itemize}
These are the only inequalities the enclosure of Section~\ref{sec:encl} ever asserts.

\section{Uniqueness as a bound on the slope of supply}\label{sec:slope}

Single crossing asks that $E=S-D$ be strictly increasing on the bracket. Since
$D'(r)\le-28.7$ on $[2.3\%,6\%]$, that holds as soon as supply does not fall faster than demand:
\begin{equation}\label{eq:target}
  S'(r)\ >\ D'(r)\quad\text{on }[2.3\%,6\%],\qquad\text{equivalently}\qquad -S'(r)\ <\ |D'(r)|.
\end{equation}
Two remarks fix what \eqref{eq:target} means and what it does not need.

First, $S$ is not differentiable everywhere. Each stage policy has a kink at the asset level where
the borrowing constraint stops binding; kinks propagate backward to wherever a policy sends a
household onto a kink of the next stage; and $S$, a finite composition of the policies along the
newborn's paths, inherits finitely many kinks in $r$. Off the kinks every policy is real-analytic in
$(a,r)$ by the analytic implicit function theorem applied to the Euler equality, so $S$ is continuous
and piecewise $C^1$ with finitely many pieces, each extending $C^1$ to its closed piece. A
continuous, piecewise-$C^1$ function with positive derivative on each open piece is strictly
increasing. So \eqref{eq:target}, read on the open pieces, suffices. (The companion's finiteness
theorem, \lean{PiecewiseAnalyticOn.finite\_zeros}, is the same observation one level down: piecewise
analytic with no flat piece gives finitely many zeros. It is not needed here, because
\eqref{eq:target} gives the stronger conclusion directly.)

Second, the true slope of supply is large and positive --- about $+53$ at $4\%$ and $+62$ at $6\%$
on the solved model --- so \eqref{eq:target} holds in fact with room to spare, and a proof of it may
be lossy by more than the true slope and still close. The construction below certifies
$S'\ge-B_d$ with $B_d$ between $6.5$ and $9.2$: a bound that does not even know supply rises, and
is still a third of what is needed.

The link from a derivative bound to the secant inequality that single crossing is usually stated
with is the mean value theorem, and it is worth writing down because it is the step that the
formalised envelope \eqref{eq:envelope} does not take.

\begin{lemma}[From the derivative to the secant]\label{lem:mvt}
Let $r_1<r_2$ lie in a cell on which, for every stage $k$, every earnings state and every asset level
in the stage's region, $c_k(a,z;\cdot)$ is continuous and piecewise $C^1$ in the rate with finitely
many pieces, and suppose
\[
  \bigl(\partial_rc_k(a,z;r)-a\bigr)^+\ \le\ \varepsilon_k(a,z)\qquad\text{at every $r$ in the cell off the kinks.}
\]
Then the failure of Theorem 1 between $r_1$ and $r_2$ at $(a,z)$ is at most
$(r_2-r_1)\,\varepsilon_k(a,z)$, and the envelope \eqref{eq:envelope} gives
\[
  \frac{S(r_2)-S(r_1)}{r_2-r_1}\ \ge\ -\frac1J\sum_kG_k\sum_z\nu(z)\max_{a\in\text{reach}}\varepsilon_k(a,z)\ =:\ -B_d .
\]
\end{lemma}

\begin{proof}
A continuous, piecewise-$C^1$ function is absolutely continuous, so
$c_k(a,z;r_2)-c_k(a,z;r_1)-(r_2-r_1)a=\int_{r_1}^{r_2}(\partial_rc_k-a)\,dr\le(r_2-r_1)\varepsilon_k(a,z)$,
and the positive part is bounded by the same. Insert $e_j=(r_2-r_1)\sum_z\nu(z)\max_a\varepsilon$
and $L_j=(1-\kappa_k)R_2$ into \eqref{eq:envelope}; the accumulation
$f_j$ is linear in the $e$'s, and $\sum_j f_j$ with the $L_j$ above is $\sum_k G_k e$ by the
definition of $G$, with $R_2$ in place of $R$. Divide by $r_2-r_1$.
\end{proof}

So the whole problem is to certify, on every cell of a partition of the bracket, a bound
$\varepsilon_k(a,z)$ on the \emph{derivative-form failure} $(\partial_rc_k-a)^+$ at every state in
the cohort's reach, and to check $B_d<\min_{\text{cell}}|D'|$. The rest of the note is that
certificate.

\section{The derivative of consumption in the rate}\label{sec:deriv}

\subsection{The formula}

Fix a stage $k$, an earnings state $z$ and an asset level $a$ at which the household saves, so that
the Euler equality holds:
\[
  c^{-\gamma}=\beta R\sum_{z'}\pi(z,z')\,c_k(x,z')^{-\gamma},\qquad x=m-c=Ra+y(z)-c,
\]
with $c=c_{k+1}(a,z)$. Differentiate in $r$ at fixed $a$, using $\partial_rR=1$ and
$\partial_rx=a-\partial_rc$:
\[
  -\gamma c^{-\gamma-1}\,\partial_rc
  =\beta\sum_{z'}\pi\,c_k^{-\gamma}
   -\gamma\beta R\sum_{z'}\pi\,c_k^{-\gamma-1}\bigl[\partial_xc_k\,(a-\partial_rc)+\partial_rc_k\bigr].
\]
Multiply by $-c^{\gamma+1}/\gamma$ and use the Euler equality on the first term,
$\beta c^{\gamma+1}\sum\pi c_k^{-\gamma}=c/R$. With the \emph{Euler weights}
\begin{equation}\label{eq:weights}
  w_{z'}=\beta R\,\pi(z,z')\,c_k(x,z')^{-\gamma-1}\,c^{\gamma+1},\qquad
  A=\sum_{z'}w_{z'},\qquad \Psi_x=\sum_{z'}w_{z'}\,\partial_xc_k(x,z'),
\end{equation}
this is $\partial_rc=-c/(\gamma R)+\Psi_x(a-\partial_rc)+\sum_{z'}w_{z'}\partial_rc_k$, hence
\begin{equation}\label{eq:dcdr}
  \partial_rc_{k+1}(a,z)
  =\frac{-\dfrac{c}{\gamma R}+\Psi_x\,a+\sum_{z'}w_{z'}\,\partial_rc_k(x,z')}{1+\Psi_x},
  \qquad
  \partial_rg_{k+1}=a-\partial_rc_{k+1}.
\end{equation}
At a corner, $c=m$ and $\partial_rc=a$ exactly. In the last period, $c_0=m$ and
$\partial_rc_0=a$ everywhere; that is the terminal condition of a backward recursion in which
\eqref{eq:dcdr} is the step.

The three ingredients of the step are bounded by things the construction controls. The weights
$w_{z'}$ are built from consumption at two consecutive stages; $A$ is the power-mean amplification
of the companion's Section 7, close to one. The slope $\partial_xc_k$ of next period's consumption in
next period's assets lies in $[\kappa_kR,R]$ by the propensity floor \eqref{eq:mpcfloor} and
monotone saving, so $\Psi_x\in[\kappa_kR\,A,\,R\,A]$. And $\partial_rc_k$ is the previous step's
output.

\subsection{The derivative-form failure and the three recursions}

The failure of Theorem 1 in derivative form is
\[
  e_{k+1}(a,z)=\bigl(\partial_rc_{k+1}(a,z)-a\bigr)^+=\bigl(-\partial_rg_{k+1}(a,z)\bigr)^+ ,
\]
the amount by which saving falls with the rate. Rewriting \eqref{eq:dcdr},
\[
  \partial_rc_{k+1}-a
  =\frac{-\dfrac{c}{\gamma R}-a+\sum_{z'}w_{z'}\,\partial_rc_k(x,z')}{1+\Psi_x},
\]
so the failure at stage $k+1$ is driven by $\sum_{z'}w_{z'}\partial_rc_k$ exceeding
$a+c/(\gamma R)$, damped by $1+\Psi_x$. Upper and lower bounds $U_k\ge\partial_rc_k\ge L_k$ propagate
through \eqref{eq:dcdr} from $U_0=L_0=a$: the numerator's upper bound uses $U_k$ where
$w_{z'}>0$, the smallest admissible $c$ and the largest admissible $a$; and the denominator's
range is the range of $\Psi_x$.

How much is kept of \eqref{eq:dcdr} decides whether the bound is useful, and the companion measured
three versions on the solved model. The \emph{crude} recursion replaces $\sum w\,\partial_rc_k$ by
$A\max_{z'}U_k$ and $\Psi_x$ by its range; it is what a proof from the two Euler inequalities alone
would deliver, and it is dead --- $B_d=312$ at $\gamma=3$ against $39$ needed, its worst state the
newborn at zero wealth. The \emph{weighted} recursion keeps $\sum_{z'}w_{z'}U_k$ with the Euler
weights and asks nothing of $\Psi_x$ beyond its range; it gives $B_d\approx4$, a margin of ten. The
\emph{exact} recursion keeps $\Psi_x$ too and is the chain rule. What separates the crude from the
weighted recursion is a single max over tomorrow's states, and that is the lesson the certificate is
built on: the weights $w_{z'}$ must be carried, which means the consumption functions at two
consecutive stages must be enclosed tightly enough that the weights are known to a few percent. The
companion's sensitivity check --- every weight, slope and interpolated bound perturbed adversarially
by one percent --- leaves the weighted bound at about $4.2$, so enclosures of that quality suffice.

\section{Enclosing the consumption functions}\label{sec:encl}

\subsection{The endogenous-grid recursion, in intervals}

The enclosure is the endogenous-grid recursion of the Euler equality, run in integers over
$D=10^{15}$ with every operation rounded outward (\texttt{numerics/rational/encl.py}, shared by
\texttt{derivcert.py}). At stage $k+1$, state $z$ and asset level $a$ with $m=y_z+Ra$, consumption
is $m$ where $\Psi_k(0)\ge m$ and otherwise the unique solution of $c=\Psi_k(m-c)$, where
\[
  \Psi_k(a')=\Bigl[\beta R\sum_{z'}\pi(z,z')\,c_k(a',z')^{-\gamma}\Bigr]^{-1/\gamma}
\]
is the consumption today that makes the Euler equality hold against saving $a'$. $\Psi_k$ is
increasing in $a'$ (monotone consumption) and increasing in $c_k$, so bracketing $c_k$ brackets
$\Psi_k$: from stage-$k$ bounds $\underline c_k\le c_k\le\overline c_k$ on the grid, build
$\Psi^{\mathrm{lo}}_k$ from $\underline c_k$ with every rounding downward and $\Psi^{\mathrm{hi}}_k$
from $\overline c_k$ with every rounding upward. Since $c\mapsto\Psi(m-c)-c$ is strictly decreasing,
\[
  c_L\le\Psi^{\mathrm{lo}}_k(m-c_L)\ \Rightarrow\ c_L\le c_{k+1}(a,z),\qquad
  c_H\ge\Psi^{\mathrm{hi}}_k(m-c_H)\ \Rightarrow\ c_H\ge c_{k+1}(a,z),
\]
which are exactly the first two clauses of the interval Euler step, with $\ell=\underline c_k$ and
$h=\overline c_k$ read at $m-c_L$ and $m-c_H$. Each certified bound costs one rigorous evaluation
of $\Psi$ from a floating-point candidate; if the candidate fails, one rigorous iteration of
$\Psi$ is itself certified (since $\Psi^{\mathrm{lo}}$ is nondecreasing, $\Psi^{\mathrm{lo}}(m-t)$
for $t\ge c^*$ is a certified lower bound), and after six iterations a bisection in rational
arithmetic is the fallback. The corner region $[0,a^*]$ on which $c=m$ exactly is certified at each
stage and state from $\Psi^{\mathrm{lo}}_k(0)$: wherever $y_z+Ra\le\Psi^{\mathrm{lo}}_k(0)$, the
third clause of the interval Euler step applies. The terminal stage is exact, $c_0=m$.

\subsection{Between grid points}

Consumption is enclosed at the grid points; the recursion needs it at the saving $m-c$, which is not
a grid point. Monotonicity and concavity close the gap: on $[a_i,a_{i+1}]$ the function lies above
the chord of the two lower bounds, and below each of three upper bounds --- the right endpoint's
upper bound (monotonicity), the left secant through the bounds at $a_{i-1},a_i$ extended forward, and
the right secant through the bounds at $a_{i+1},a_{i+2}$ extended backward --- of which the smallest
is taken. Above the certified corner the chord is anchored at the corner itself, where consumption
is known exactly. These are the last clause of the interval Euler step.

\subsection{The grid, and the widths}

The grid has $1200$ points from $0$ to $3$ in steps of $0.0025$ and then relative steps of
$0.25\%$ up to $3000$: $3968$ points in all. The top is far above the reach of any cohort (about
$44$) because unavailability beyond the grid's top creeps inward by a factor of about $1/R$ per
stage and must never reach the cohort's region; the fine spacing is for the derivative recursion,
whose interpolation slack is of order $c''h^2$ per stage accumulated over a transient of some fifty
stages. Two details of the implementation mattered for the widths. The floating-point candidate for
each certified fixed point must be computed on its own side of the enclosure --- the lower
candidate from $\Psi^{\mathrm{lo}}$, the upper from $\Psi^{\mathrm{hi}}$ --- or each stage loses a
fraction $\Psi_x$ of the width; and the corner test must be strict (the corner at zero wealth is
certified only where the household demonstrably does not save there). With both, the enclosure
width at the newborn stage over the cohort's region is $0.0023$ at $4\%$ and $0.0015$ at $6\%$,
against consumption of order one, which is what the weights need.

What the enclosure delivers at a single rate is, at every stage $k$, every state $z$ and every
grid point $a_i$, certified numbers $\underline c_k(a_i,z)\le c_k(a_i,z)\le\overline c_k(a_i,z)$ and
a certified corner $a^*_k(z)$, each one an instance of the interval Euler step; and from these, by
the chord and secants, bounds at every asset level up to the grid's top. The cohort's reach
$\rho_j$ is computed from the same data: $\rho_0=0$ and $\rho_{j+1}=\max_z\bigl(m(\rho_j)-
\underline c_{J-1-j}(\rho_j,z)\bigr)$, the most any household can hold at age $j+1$, since saving is
nondecreasing in assets.

\section{The interval recursion for the derivative}\label{sec:recursion}

A derivative has no monotonicity or concavity to interpolate with, so the bounds on $\partial_rc_k$
are carried on grid \emph{intervals} rather than at grid points: $U_k[z][i]\ge\partial_rc_k(a,z)\ge
L_k[z][i]$ for every $a\in[a_i,a_{i+1}]$. The terminal values are $U_0[z][i]=a_{i+1}$ and
$L_0[z][i]=a_i$, since $\partial_rc_0=a$. One step, from stage $k$ to stage $k+1$ on the interval
$[a_i,a_{i+1}]$ at state $z$, evaluates \eqref{eq:dcdr} with every ingredient replaced by its range
over the interval and rounded outward:
\begin{itemize}[leftmargin=1.5em,itemsep=2pt]
\item today's consumption $c\in[\underline c_{k+1}(a_i,z),\overline c_{k+1}(a_{i+1},z)]$, by
  monotonicity;
\item today's saving $x=m-c\in[m(a_i)-\overline c_{k+1}(a_i,z),\ m(a_{i+1})-\underline c_{k+1}(a_{i+1},z)]$,
  since saving is nondecreasing in assets;
\item tomorrow's consumption $c_k(x,z')$ in the enclosure at the two ends of that $x$-range, by
  monotonicity;
\item each weight $w_{z'}$ between its two extreme values over those ranges, and $\Psi_x$ in
  $[\kappa_kR\,A,\,R\,A]$ with $A=\sum w$ and $\kappa_k$ from its certified lower bound;
\item $\partial_rc_k(x,z')$ between the smallest $L_k$ and the largest $U_k$ over the grid intervals
  that the $x$-range meets.
\end{itemize}
The interior formula needs the Euler equality, so it is used only where saving is positive on the
whole interval, $m(a_i)-\overline c_{k+1}(a_i,z)>0$. On every other interval --- those inside the
certified corner, where $\partial_rc=a$ exactly, and those that straddle the point where the
constraint stops binding, where the household may be at the corner at some $a$ in the interval ---
the bounds are widened to cover $\partial_rc=a$ as well: $U\leftarrow\max(U,a_{i+1})$,
$L\leftarrow\min(L,a_i)$. Beyond the grid's top no bound is available, and that unavailability
creeps inward from $3000$ by about one interval per stage, never approaching the reach.

The failure bound on the interval is $\varepsilon_{k+1}[z][i]=(U_{k+1}[z][i]-a_i)^+$, the largest
value $\partial_rc-a$ can take there. By Lemma~\ref{lem:mvt} what enters the envelope is, for each
age $j$ (stage $J-j$), $\sum_z\nu(z)\max\{\varepsilon_{J-j}[z][i]:a_i\le\rho_{j-1}\}$, the maximum
over the intervals inside the reach; and then
\[
  B_d=\frac1J\sum_{j=1}^{J}f_j,\qquad f_1=0,\quad
  f_{j+1}=(1-\kappa_{J-j})R\,f_j+\sum_z\nu(z)\max_{\text{reach}}\varepsilon_{J-j}[z][\cdot],
\]
with $\kappa$ its certified lower bound (which makes $1-\kappa$ an upper bound) and $R$ the gross
rate. Every number in this display is a rational, computed exactly.

At a single rate on the fine grid this gives, for $\gamma=3$, $B_d=5.2$ at $4\%$ against
$|D'|=39.1$ and $B_d=7.4$ at $6\%$ against $28.7$; and for $\gamma=5$, $B_d=12.8$ at $4\%$. The
failure is concentrated at the youngest ages --- about $0.4$ per age at ages one to ten, $0.23$ at
age thirty --- which is where the enclosures are widest and where, on the solved model, the true
derivative-form failure is exactly zero at $\gamma=3$: the bound is interpolation slack, not
economics. That is why the grid is fine.

\section{From a rate to a cell of rates}\label{sec:cell}

Lemma~\ref{lem:mvt} needs the derivative-form failure bounded at every rate in the cell, so the
enclosures the recursion reads must hold at every rate in the cell. A direct cell enclosure --- the
same endogenous-grid recursion with the two rates at the ends of the cell placed where each hurts
--- is valid but too loose: its widths are twenty to seventy times the single-rate widths on cells
of a tenth of a percentage point, and the recursion's weights carry consumption to the power
$\gamma+1$. A cell enclosure fed back from the recursion's own derivative bound diverges for the
same reason. What works is to \emph{freeze} the transport and verify it afterwards.

\subsection{The frozen transport}

Let $h$ be the cell's width and $r_0$ its midpoint. Run the point enclosure at $r_0$ to get
$[\underline c^0_k,\overline c^0_k]$ at every stage, state and grid point, and run the derivative
recursion at $r_0$ to get $U^0_k$. Now \emph{assume} that on the cell, at every stage and every
state on a wealth zone $Z_k$ to be specified,
\begin{equation}\label{eq:assume}
  |\partial_rc_k(a,z;r)|\ \le\ M_k[z][i]:=\lambda\,\max\bigl(|U^0_k[z][i]|,|L^0_k[z][i]|\bigr)+\mu,
  \qquad a\in[a_i,a_{i+1}]\subset Z_k,
\end{equation}
with $\lambda=2$ and $\mu=1$. Under the assumption, consumption at any rate in the cell is within
$M_kh/2$ of its value at the midpoint, so the transported enclosure
\[
  \underline c_k(a_i,z)=\underline c^0_k(a_i,z)-M_kh/2,\qquad
  \overline c_k(a_i,z)=\overline c^0_k(a_i,z)+M_kh/2
\]
holds on the whole cell on the zone. Outside the zone the enclosure falls back to bounds that need
no assumption at all: the propensity floor $\kappa_km$ from below, read at the lower rate, and cash
on hand $m$ from above, read at the upper rate. Then run the derivative recursion on the transported
enclosures, with every rate-dependent quantity taken at whichever end of the cell makes the bound
larger, to get bounds $U_k$ valid at every rate in the cell \emph{at which the transported
enclosure holds}, and check, on every interval where \eqref{eq:assume} was assumed,
\begin{equation}\label{eq:verify}
  \max\bigl(|U_k[z][i]|,|L_k[z][i]|\bigr)\ <\ M_k[z][i]\qquad\text{strictly.}
\end{equation}

\begin{lemma}[Continuation]\label{lem:cont}
Suppose that for every stage $k$, state $z$ and asset level $a$ in the cohort's region,
$r\mapsto c_k(a,z;r)$ is continuous and piecewise $C^1$ on the cell with finitely many pieces, the
derivative extending continuously to each closed piece; and that the derivative recursion is sound:
whenever the transported enclosure holds at a rate $r$ in the cell on the zone, and the fallback
bounds outside it, the recursion's output satisfies $L_k\le\partial_rc_k(\cdot;r)\le U_k$ on every
interval, one-sidedly at kinks. If \eqref{eq:verify} holds wherever \eqref{eq:assume} was assumed,
then \eqref{eq:assume} is true on the whole cell, and so the transported enclosure and the bounds
$U_k,L_k$ hold at every rate in the cell.
\end{lemma}

\begin{proof}
Let $T$ be the set of $t\in[0,h/2]$ such that $|c_k(a,z;r)-c_k(a,z;r_0)|\le M_k|r-r_0|$ for all
$|r-r_0|\le t$, all stages, states and asset levels in the zone. $T$ contains $0$, and is closed
because each inequality is non-strict and $c_k$ is continuous in $r$. Let $t\in T$ with $t<h/2$. For
every $r$ with $|r-r_0|\le t$ the transported enclosure holds (the midpoint enclosure holds at
$r_0$, and consumption has moved by at most $M_kt\le M_kh/2$), so by soundness
$|\partial_rc_k|\le\max(|U_k|,|L_k|)<M_k$ on the zone at $r$, one-sidedly at kinks. The one-sided
derivatives at $r_0\pm t$ are strictly below $M_k$, and the derivative extends continuously to the
closed pieces, so there is $\varepsilon>0$ with $|\partial_rc_k|<M_k$ on
$|r-r_0|<t+\varepsilon$ off the kinks; by absolute continuity
$|c_k(r)-c_k(r_0)|\le\int|\partial_rc_k|\le M_k|r-r_0|$ there, so $t+\varepsilon/2\in T$. A closed
subset of $[0,h/2]$ that contains $0$ and has no largest element below $h/2$ is all of $[0,h/2]$.
\end{proof}

The hypotheses of the lemma are the regularity of Section~\ref{sec:slope} and the soundness of the
recursion, which is Section~\ref{sec:recursion} read with the cell's rate range in place of a rate.
The lemma is the one piece of the construction that is neither a formalised theorem nor arithmetic:
it is a paper argument, and a short one.

\subsection{What makes it run}

Three refinements were needed, and each failed version is recorded in the companion.

\emph{The zone.} The assumption \eqref{eq:assume} cannot be made on the whole grid: near the grid's
top the midpoint bounds $U^0$ are enormous and the transported widths with them, and a few grid
points just below the top then fail \eqref{eq:verify}. The assumption is confined to the zone the
recursion actually reads from the cohort's region, $Z_{J-1}=[0,1.2\rho+2]$ at the newborn stage and
$Z_k=[0,\max(1.2\rho_{\mathrm{age}}+2,\ (1-\kappa_{k+1})(y_{\max}+R\sup Z_{k+1})+2)]$ working
backward, which is closed under the saving map by the propensity floor; its top runs from about
$32$ at the newborn stage through $137$ at stage thirty to $2$ at the last stage. The recursion
skips intervals above the zone, where the fallback bounds would make the rationals enormous without
being read.

\emph{The floor $\mu$.} Near the corner the midpoint bound $U^0$ is of order $a$, tiny, and the
cell pass --- which cannot see that the household is constrained at every rate in the cell ---
returns a bound above $\lambda U^0$. The additive floor $\mu=1$ covers those states.

\emph{The multiplier $\lambda$ does not help.} Cells of $0.02$ percentage points fail
\eqref{eq:verify} by nine percent at every $\lambda$ tried, because the transported width scales
with $\lambda$ and so does the bound it produces; cells of $0.1$ points fail it forty times over.
Cells of $0.01$ points verify at a ratio of $0.75$, and cells of $0.005$ at $0.63$. The width is
what it is: $0.01$ percentage points at the present enclosure quality.

The cost of the cell is the \emph{inflation} of $B_d$ over the midpoint's: at $5\%$ the cell
bound is $6.9$ against a midpoint bound of $6.3$, nine percent; at $6\%$, $8.1$ against $7.4$. The
transport does not spoil the margin.

\section{Covering the bracket}\label{sec:bracket}

The bracket $[2.3\%,6\%]$ is partitioned into $370$ cells of $0.01$ percentage points, run in six
batches: $[4\%,6\%]$ first, then five extensions downward to $2.3\%$, where the Lean theorem of
Section~\ref{sec:outer} takes over. For each cell the certificate runs the point enclosure at the
midpoint (about twenty-four minutes on the fine grid), the derivative recursion at the midpoint, the
transported enclosure and the cell recursion, and the verification \eqref{eq:verify}; a cell that
failed would be split in two and both halves rerun, down to a minimum width of $0.001$ points. Each
batch is resumable, with one line per finished cell in
\texttt{numerics/rational/runs/bracket\_g3\_*.txt}; the whole took about $470$ processor-hours on
seven cores, at $76$ minutes per cell.

No cell needed splitting. Every cell certified at its first width, with the verification ratio
$\max|U|/M$ about three quarters on every one of them, and the certified bounds are:

\begin{center}
\begin{tabular}{@{}lcccc@{}}
\toprule
sub-range & $\max B_d$ & $\min|D'|$ & worst $B_d/|D'|$ & cells\\
\midrule
$[2.3\%,3.0\%]$ & $5.10$ & $46.49$ & $0.110$ & $70$\\
$[3.0\%,4.0\%]$ & $6.45$ & $39.06$ & $0.165$ & $100$\\
$[4.0\%,5.0\%]$ & $7.82$ & $33.28$ & $0.235$ & $100$\\
$[5.0\%,6.0\%]$ & $9.15$ & $28.70$ & $0.319$ & $100$\\
\bottomrule
\end{tabular}
\end{center}

$B_d$ rises smoothly from $4.28$ at the bottom of the bracket to $9.15$ at the top, while $|D'|$
falls, so the worst ratio, $0.319$, sits at the top cell $[5.99\%,6.00\%]$. On every cell
$B_d<\min_{\text{cell}}|D'|$, with a margin of at least three. By Lemma~\ref{lem:mvt} applied cell
by cell, for any two rates $r_1<r_2$ in the same cell,
\[
  E(r_2)-E(r_1)\ \ge\ (r_2-r_1)\bigl(\min\nolimits_{\text{cell}}|D'|-B_d\bigr)\ >\ 0,
\]
and since adjacent cells share an endpoint, $E$ is strictly increasing on the whole bracket.

\section{The bracket contains an equilibrium}\label{sec:level}

Strict monotonicity of $E$ gives at most one zero. For exactly one, $E$ must change sign on the
bracket, and the same enclosures certify that it does (\texttt{numerics/rational/levelcert.py}).
Run the cohort-assets recursion of Section~\ref{sec:economy} with the \emph{interval} policy. At a
grid point $a_i$ and state $z$ the saving $g_{k+1}(a_i,z)$ lies in
$[m(a_i)-\overline c_{k+1}(a_i,z),\ m(a_i)-\underline c_{k+1}(a_i,z)]^+$ by the enclosure, with $m(a_i)$ rounded down in the first and up in the second; $A_k$ is
nondecreasing in assets, because saving is and the recursion preserves monotone functions; so a lower
bound on $A_{k+1}(a_i,z)$ is $a_i$ plus the $\pi$-expectation of the lower bound on $A_k$ read at the
grid point \emph{below} the lower bound on saving, and an upper bound is the same with upper bounds
and the grid point \emph{above}, each rounded outward. Sixty steps from $A_0=a$ give
$S(r)\in[\underline S,\overline S]$:

\begin{center}
\begin{tabular}{@{}lccc@{}}
\toprule
$r$ & $S(r)$ certified & $D(r)$ & verdict\\
\midrule
$4\%$ & $[3.821,\,3.966]$ & $4.6875$ & $S<D$\\
$6\%$ & $[4.950,\,5.159]$ & $4.0179$ & $S>D$\\
\bottomrule
\end{tabular}
\end{center}

The widths are the grid rounding of sixty compositions and are of no consequence against gaps of
$0.72$ and $0.93$. So $E(4\%)<0<E(6\%)$.

\section{Proof of Theorem~\ref{thm:main}}\label{sec:proof}

$S$ is continuous on $[2.3\%,6\%]$ (\lean{continuousOn\_olgSupply}); the hypotheses of that theorem
are the stage properties on a reachable family, and the enclosure's reach, $\rho_j\le44$, is such a
family for any cap above it, with cap slack certified by the enclosure itself (the upper bound on
saving never approaches the grid's top). $D$ is continuous. By Section~\ref{sec:level},
$S(4\%)<D(4\%)$ and $S(6\%)>D(6\%)$, so by the intermediate value theorem
(\lean{exists\_olgEquilibrium}) there is a rate in $(4\%,6\%)$ with $S=D$. By
Section~\ref{sec:bracket}, $E=S-D$ is strictly increasing on $[2.3\%,6\%]$, so that rate is the
only one in $[2.3\%,6\%]$. By Theorem~\ref{thm:outer} of Section~\ref{sec:outer}, $S<D$ at every
rate in $(-8\%,2.3\%]$, so there is none there. \qed

Read backwards, the proof uses: the certified endpoint levels (Section~\ref{sec:level}, arithmetic
on the enclosures); the certified slope bound on every cell (Sections~\ref{sec:recursion}
and~\ref{sec:cell}, arithmetic on the enclosures plus Lemma~\ref{lem:cont}); the enclosures
themselves (Section~\ref{sec:encl}, instances of the interval Euler step); Lemma~\ref{lem:mvt} and
the derivative formula \eqref{eq:dcdr} (paper); and the companion's theorems of
Section~\ref{sec:known} (kernel-checked).

\section{Outside the bracket, by proof term}\label{sec:outer}

Below and above the bracket the question changes character. Near the equilibrium only the slope
certificate can separate supply from demand, because the proved level bounds are too loose there;
away from it the level bounds of the companion suffice, and they are cheap enough to evaluate inside
Lean on whole ranges of rates.

\paragraph{The two bounds.} The companion proves an affine ceiling on accumulated assets,
$A_k(a,z)\le\bar F_k(z)+G_ka$ with $\bar F$ an explicit recursion in $\kappa_k$, the cohort slope
$G_k$ and a risk-adjusted human wealth built from power means of exponent $-\gamma$, valid whenever
the cap is slack everywhere (\lean{olgCapital\_le\_cohortCeil}); and a clamped floor, the recursion
$\Lambda_0(z,a)=a$, $\Lambda_{k+1}(z,a)=a+\sum_{z'}\pi(z,z')\Lambda_k(z',\Phi_{k+1}(z,a))$ with
$\Phi$ the saving floor $\max\{0,(1-\kappa_k)(y_z+Ra)-\kappa_kH_k(z)\}$, valid whenever a lifetime
of feasible saving stays under the cap (\lean{le\_olgCapital\_at}). Supply is at most the
$\nu$-average of $\bar F_{59}$ and at least that of $\Lambda_{59}(\cdot,0)$, each over sixty.

\paragraph{From a rate to a cell.} Every ingredient is monotone in the rate $R$ or in the patience
factor $\Thorn=(\beta R)^{1/3}$: $\kappa_k$ rises with $R$ and falls with $\Thorn$, human wealth
falls with $R$, the cohort slope rises with $R$ and falls with $\kappa$, and the power mean is
monotone in its arguments. On a cell $[r_1,r_2]$ a checker (\texttt{EnclCert/Outer.lean})
brackets $\Thorn$ by rational cube-root bounds, runs each recursion with every ingredient taken at
the end of the cell that makes the bound larger, rounding outward over $D=10^{12}$, and compares
the ceiling with demand at $r_2$ or the floor with demand at $r_1$. The floor is carried on a grid of
$650$ asset levels and read at the grid point below the certified saving floor, which is sound
because $\Lambda_k$ is monotone in assets. Every stored value is re-verified by an exact integer
inequality with denominators cleared, and only those inequalities are used by the bridge
(\texttt{OLG/OuterBridge.lean}), which proves by induction on the stage that each stored number
bounds the real quantity it stands for, discharges the slackness hypothesis of the ceiling from
impatience at every stage (the limiting propensity $1-\Thorn/R$ where $\Thorn<R$, and
$y_{\max}+r\bar a\le0$ at rates negative enough), and checks the reach hypothesis of the floor at
the top of the cell. The cap is $10^6$, which the reach at $20\%$ needs; the ceiling side is
indifferent to it.

\paragraph{The meshes.} The ceiling crosses demand near $2.35\%$ and the floor near $8.4\%$ at this
calibration. The low mesh is $165$ cells from $-8\%$ to $2.3\%$, of $0.1$ points to $1.5\%$,
$0.02$ to $2.1\%$ and $0.005$ to $2.3\%$; the high mesh is $121$ cells from $8.7\%$ to $20\%$, of
$0.01$ points to $9.5\%$, $0.1$ to $12\%$ and $0.5$ to $20\%$. Both verdicts are established by
\texttt{native\_decide} (\lean{lowCheck\_eq\_true}, \lean{highCheck\_eq\_true}), about four
minutes of compiled arithmetic each; a covering lemma puts every rate of a range in some cell of an
increasing mesh.

\begin{theorem}[{No equilibrium outside $[2.3\%,8.7\%]$}]\label{thm:outer}
At $\gamma=3$ and the calibration of Section~\ref{sec:economy} with cap $10^6$, capital supply is
strictly below demand at every rate in $(-8\%,2.3\%]$ and strictly above it at every rate in
$[8.7\%,20\%]$.
\formal{EnclCert.olg7\_supply\_lt\_demand\_low, EnclCert.olg7\_demand\_lt\_supply\_high}
\end{theorem}

Unlike Theorem~\ref{thm:main}, this one is a Lean theorem: its only axiom beyond classical Mathlib
is the one \texttt{native\_decide} introduces for the two compiled verdicts. What it leaves
uncovered below $20\%$ is $(6\%,8.7\%]$, where neither bound separates supply from demand and only
the slope certificate of Sections~\ref{sec:recursion}--\ref{sec:bracket} can; extending that
certificate to it is about $270$ further cells of the same run. Above
$20\%$ the formal model's cap starts to bind for a household that saves everything feasible, and
the capped economy stops being the one of interest.

\section{Standing, and what remains}\label{sec:standing}

\subsection{What kind of result this is}

\begin{center}
\begin{tabular}{@{}p{0.47\textwidth}p{0.47\textwidth}@{}}
\toprule
step & standing\\
\midrule
stage properties, propensity floor, cohort slope, earnings marginal, cohort envelope, continuity,
existence, interval Euler step, chord and secants & theorems, checked by the Lean kernel\\[2pt]
the derivative formula \eqref{eq:dcdr}; Lemma~\ref{lem:mvt}; Lemma~\ref{lem:cont};
piecewise-$C^1$ regularity of the policies in the rate & paper arguments, given above\\[2pt]
the enclosures at two hundred midpoints and two endpoints; the derivative recursion on each cell;
the verification \eqref{eq:verify}; the endpoint levels & exact rational arithmetic, rounded
outward, evaluated by \texttt{derivcert.py} and \texttt{levelcert.py}\\
\bottomrule
\end{tabular}
\end{center}

The arithmetic row is where the standing is weakest, and it is weak in a specific way: each
inequality the programs check is an instance of a formalised theorem (the enclosure) or of a
recursion written out in Sections~\ref{sec:recursion}--\ref{sec:level} (the derivative bound, the
level bound), but the programs that evaluate them were not proved to evaluate what this note says
they evaluate. The companion has already removed that gap once, for the pairwise enclosure
certificate at $(4\%,4.6\%)$: a Lean checker over naturals, its verdict established by
\texttt{native\_decide} in a quarter of an hour, and a bridge from the checker's inequalities to the
theorems behind them (\lean{enclCheck\_eq\_true}, \lean{olg7\_increment}). The derivative certificate
has the same shape, with three additions: the interval recursion of Section~\ref{sec:recursion} as a
checker, the level recursion of Section~\ref{sec:level}, and Lemma~\ref{lem:cont} as a theorem, the
one genuinely new piece of mathematics. At seventy-six minutes per cell in Python, three hundred and
seventy cells in compiled Lean is a weekend of machine time, not a new argument.

\subsection{What the theorem does not say}

It is one calibration: one chain, one $\beta$, one $\gamma$, one technology. Nothing transfers to a
neighbouring calibration except by rerunning the certificate, which is the price of a certificate
over a proof. It is one bracket, nearly: Theorem~\ref{thm:outer} excludes every rate in $(-8\%,2.3\%]$ and in
$[8.7\%,20\%]$ by proof term, so what remains open below $20\%$ is $[2.3\%,4\%)$ and
$(6\%,8.7\%]$, where the slope certificate has not yet been run; until it is, the honest statement
is exactly one equilibrium on $(-8\%,6\%]$ and none in $[8.7\%,20\%]$. And it is $\gamma=3$: the same run at $\gamma=5$, where the
single-rate bound is $12.8$ against $39$ at $4\%$, is the same command with one argument changed and
about the same machine time, but it has not been run.

\subsection{What would raise it}

In order of return on effort. A \emph{transported} derivative bound in place of the constant
multiple $\lambda U^0$ in \eqref{eq:assume} --- carrying the assumption through the recursion
rather than freezing it --- should admit cells three to five times wider, since the verification
ratio is pinned at $0.75$ by the constant multiple and not by the margin. The $\gamma=5$ bracket,
which turns the theorem into a statement about Aiyagari's second calibration. The Lean replay of the
derivative checker with Lemma~\ref{lem:cont}, which turns the certificate into a theorem with one
\texttt{native\_decide} axiom. And the formalisation of piecewise analyticity of the stage policies
in the rate, which the regularity hypotheses of Lemmas~\ref{lem:mvt} and~\ref{lem:cont} rest on and
which is at present a paper argument from the analytic implicit function theorem.

\appendix
\section{Reproduction}

All files are under \texttt{WriteUps/WriteUpOLG/numerics/rational/} in the repository.
\begin{itemize}[leftmargin=1.5em,itemsep=2pt]
\item The bracket runs (resumable, results appended to the file named), with the environment\\
  \texttt{ENCL\_FINE=fine ENCL\_NOCACHE=1 ENCL\_MU=1 ENCL\_LAMS=2 ENCL\_PROCS=7}\\
  \texttt{ENCL\_RESULTS=runs/bracket\_g3\_<name>.txt} and the arguments
  \texttt{derivcert.py 3 r\_lo r\_hi n 1 fixed}: \texttt{0.04 0.06 200} (4to6), \texttt{0.037 0.04 30}
  (37to40), \texttt{0.031 0.037 60} (31to37), \texttt{0.0275 0.031 35} (275to31),
  \texttt{0.0233 0.0275 42} (233to275), \texttt{0.023 0.0233 3} (230to233).
\item A batch's summary (leaves, splits, worst margin, coverage): \texttt{python3
  runs/bracket\_summary.py runs/bracket\_g3\_4to6.txt 0.04 0.06}.
\item A single rate, for the first table of Section~\ref{sec:recursion}:
  \texttt{ENCL\_FINE=fine python3 derivcert.py 3 0.04 0.04 0}.
\item The endpoint levels: \texttt{ENCL\_FINE=fine python3 levelcert.py 3 1/25} and
  \texttt{... 3 3/50}.
\item The derivative-form measurements on the solved model, which motivated the weighted recursion:
  \texttt{numerics/slopebnd.m}, \texttt{slopebnd2.m}, \texttt{dcdr.m}.
\end{itemize}
The calibration is in the header of \texttt{derivcert.py}: $J=60$, $n_z=7$, $\beta=24/25$,
$\rho=9/10$, $\alpha=9/25$, $\delta=2/25$, the Rouwenhorst chain, the seven earnings levels and
their binomial invariant law, the grid, and $D=10^{15}$.

\bibliographystyle{ecta}
\nocite{Aiyagari1994,Light2020,CarrollKimball1996,Huggett1996}
\bibliography{refs}

\end{document}
